\section{Funktionen}
\subsection{Default Argumente}
\begin{itemize}
    \item \textbf{Prototipo}: I valori standard vengono assegnati \textbf{solo} nell'interfaccia.
    \item \textbf{Omissione}: I parametri di default possono essere omessi durante la chiamata.
    \item \textbf{Ordine}: I parametri di default devono stare \textbf{\cor{sempre alla fine (a destra)}} della lista dei parametri.\end{itemize}

\lstinputlisting[language=c++]{code/02_funktionen/default_argumente.cpp}

Default Argumente können überall verwendet werden bis auf Aufrufparameter von Operatoroverloading. 
Dort sind sie \textbf{nicht} erlaubt. 
Dort machen sie aber ohnehin keinen Sinn. 

\subsection{Funzioni Inline}
\begin{itemize}
    \item \textbf{Sostituzione dei macro}: Alternativa sicura alle macro C.
    \item \textbf{Zero Overhead}: Il codice viene inserito direttamente (nessun salto di funzione).
    \item \textbf{Sicurezza dei tipi}: Il compilatore garantisce la verifica dei tipi.
    \item \textbf{Obiettivo}: Funzioni piccole, richiamate molto frequentemente (ad es. loop).
    \item \textbf{Divieto}: Non applicabile in caso di ricorsione o puntatori a funzioni.
    \item \textbf{Solo header}: Se usato negli header, la definizione della Fkt. direttamente nel file.h
    \item \textbf{Ottimizzazione}: Richiede flag attivi (ad es. \texttt{-O3}).\end{itemize}

\lstinputlisting[language=C++]{code/02_funktionen/02_funktionen_snippet_1.cpp}